\chapter[Event Sourcing y CQRS]{\emj{📨}~Event Sourcing y CQRS}

\section[Event Sourcing: estado como historia]{\emj{📜}~Event Sourcing: estado como historia}

En un modelo tradicional, se guarda el \textbf{estado actual}.
En Event Sourcing, se guarda la \textbf{historia de eventos}; el estado actual se deriva de ellos.

\begin{teoriabox}{Comparativa: estado vs eventos}
\begin{itemize}
  \item \textbf{CRUD tradicional:} \texttt{UPDATE orders SET status='shipped'} — el estado anterior se pierde
  \item \textbf{Event Sourcing:} se inserta un evento \texttt{\{type: OrderShipped, orderId: 42, at: ...\}} — el estado anterior está implícito en la historia
  \item \textbf{Ventajas:} audit log gratuito, time travel (reconstruir estado en cualquier momento), debugging trivial, event replay para nuevas proyecciones
  \item \textbf{Desventajas:} complejidad mayor, eventual consistency en las vistas (read models), el estado actual requiere replay de todos los eventos (mitigado con snapshots)
\end{itemize}
\end{teoriabox}

\section[Event Store con PostgreSQL]{\emj{🗄️}~Event Store con PostgreSQL}

PostgreSQL es una base excelente para un event store: append-only, transaccional, con índices eficientes.

Un event store bien diseñado en PostgreSQL aprovecha sus garantías ACID para garantizar que los eventos se escriben en orden y sin duplicados.
La restricción \texttt{UNIQUE (stream\_id, stream\_version)} es el mecanismo central: si dos transacciones concurrentes intentan escribir la versión 5 del mismo stream, solo una tendrá éxito.
Esta es la implementación de \textbf{Optimistic Concurrency Control}: no se bloquea nada al leer, pero al escribir se verifica que nadie más escribió primero.

\begin{lstlisting}[caption={Schema de Event Store}]
-- Tabla principal de eventos:
CREATE TABLE events (
  event_id        UUID        NOT NULL DEFAULT gen_random_uuid(),
  stream_id       TEXT        NOT NULL,   -- ej: 'order-42', 'user-100'
  stream_version  BIGINT      NOT NULL,   -- versión del stream (para OCC)
  event_type      TEXT        NOT NULL,   -- 'OrderCreated', 'OrderShipped', etc.
  payload         JSONB       NOT NULL,
  metadata        JSONB       NOT NULL DEFAULT '{}',
  created_at      TIMESTAMPTZ NOT NULL DEFAULT NOW(),

  PRIMARY KEY (event_id),
  UNIQUE (stream_id, stream_version)  -- evita conflictos de concurrencia
);

CREATE INDEX idx_events_stream ON events (stream_id, stream_version);
CREATE INDEX idx_events_type ON events (event_type);
CREATE INDEX idx_events_created_at ON events (created_at);

-- Tabla de snapshots (para reconstrucción rápida):
CREATE TABLE snapshots (
  stream_id       TEXT        NOT NULL,
  stream_version  BIGINT      NOT NULL,
  state           JSONB       NOT NULL,
  created_at      TIMESTAMPTZ NOT NULL DEFAULT NOW(),
  PRIMARY KEY (stream_id, stream_version)
);
\end{lstlisting}

\begin{lstlisting}[caption={Append de eventos con Optimistic Concurrency Control}]
-- Agregar evento verificando versión esperada (OCC):
-- Si alguien más modificó el stream, la inserción falla por UNIQUE constraint
INSERT INTO events (stream_id, stream_version, event_type, payload)
VALUES (
  'order-42',
  5,   -- esperamos que la versión actual sea 4, insertamos la 5
  'OrderShipped',
  '{"carrier": "DHL", "tracking_code": "1Z999AA1"}'::jsonb
);
-- Si ya existe (stream_id, stream_version=5): ERROR duplicate key
-- La app debe hacer retry o reportar conflict

-- Obtener versión actual de un stream:
SELECT MAX(stream_version) FROM events WHERE stream_id = 'order-42';

-- Reconstruir estado de un aggregate (replay):
SELECT event_type, payload, created_at
FROM events
WHERE stream_id = 'order-42'
ORDER BY stream_version ASC;
-- La app aplica cada evento al estado inicial vacío para obtener el estado actual
\end{lstlisting}

\section[CQRS: separar lectura de escritura]{\emj{✂️}~CQRS: separar lectura de escritura}

El event store es excelente para escribir (append-only, siempre rápido) pero muy ineficiente para leer estado actual: reconstruir el estado de una orden requiere reproducir 50 eventos cada vez.
CQRS resuelve esto separando completamente el modelo de escritura del modelo de lectura.

Las \textbf{proyecciones} (o read models) son vistas desnormalizadas, específicamente diseñadas para los patrones de consulta de la UI.
Un worker de proyección escucha el stream de eventos y actualiza estas vistas inmediatamente después de cada evento.
El resultado es eventual consistency: por breves milisegundos, la UI puede mostrar datos levemente desactualizados, pero el sistema gana velocidad de escritura ilimitada y lecturas instantáneas.

CQRS (Command Query Responsibility Segregation) separa el modelo de escritura (Commands) del modelo de lectura (Queries/Projections).

\begin{teoriabox}{CQRS con Event Store}
\begin{itemize}
  \item \textbf{Write side:} recibe Commands (\texttt{ShipOrder}, \texttt{CancelOrder}), valida, aplica lógica de negocio, emite Events al event store
  \item \textbf{Read side (Projections):} procesa events del store y actualiza vistas desnormalizadas optimizadas para lectura
  \item \textbf{Eventual consistency:} hay un delay entre que un evento se escribe y que la proyección lo refleja (generalmente milisegundos)
  \item Las proyecciones se pueden reconstruir completamente replaying todos los eventos desde el inicio
\end{itemize}
\end{teoriabox}

La tabla de proyección se diseña exclusivamente para las necesidades de la interfaz de usuario, no para reflejar la estructura del dominio.
Si la UI necesita mostrar nombre del cliente, total de la orden e ícono de estado en una lista, la proyección tiene exactamente esas columnas, ya calculadas, listas para servir en un solo SELECT sin joins.
El campo \texttt{stream\_version} actúa como checkpoint: el worker de proyección sabe hasta qué evento ha procesado y puede retomar desde ahí si reinicia.

\begin{lstlisting}[caption={Tabla de proyección (read model)}]
-- Projection optimizada para UI de órdenes:
CREATE TABLE order_summary (
  order_id        BIGINT      PRIMARY KEY,
  customer_name   TEXT        NOT NULL,
  status          TEXT        NOT NULL,
  total_amount    NUMERIC     NOT NULL,
  item_count      INT         NOT NULL DEFAULT 0,
  last_event_at   TIMESTAMPTZ,
  stream_version  BIGINT      NOT NULL  -- checkpoint del último evento procesado
);

-- Actualizar projection al procesar un evento:
-- (esto corre en el projection worker, no en el request path)
INSERT INTO order_summary (order_id, customer_name, status, total_amount, stream_version)
VALUES (42, 'Ana García', 'created', 0, 1)
ON CONFLICT (order_id) DO UPDATE
  SET status = EXCLUDED.status,
      stream_version = EXCLUDED.stream_version,
      last_event_at = NOW();
\end{lstlisting}

\section[Outbox Pattern]{\emj{📤}~Outbox Pattern}

Problema: guardar un evento en la BD Y publicarlo en Kafka en la misma operación. Si el mensaje a Kafka falla después de que la BD committeó, hay inconsistencia.

\begin{teoriabox}{Solución: Outbox}
\begin{enumerate}
  \item La transacción de negocio escribe el evento en la BD \textbf{y} en una tabla \texttt{outbox} — todo en la misma transacción ACID.
  \item Un proceso separado (poller o CDC) lee la tabla outbox y publica los mensajes a Kafka.
  \item Marca los mensajes como procesados (o los elimina).
  \item Si el publisher falla, los mensajes siguen en el outbox y se reintentarán. \textbf{At-least-once delivery}.
\end{enumerate}
\end{teoriabox}

\begin{lstlisting}[caption={Outbox table y proceso}]
-- Tabla outbox:
CREATE TABLE outbox (
  id          UUID        NOT NULL DEFAULT gen_random_uuid() PRIMARY KEY,
  topic       TEXT        NOT NULL,      -- 'orders.events', 'payments.events'
  key         TEXT,                      -- para particionamiento en Kafka
  payload     JSONB       NOT NULL,
  created_at  TIMESTAMPTZ NOT NULL DEFAULT NOW(),
  sent_at     TIMESTAMPTZ              -- NULL = pendiente
);

-- En la transacción de negocio:
BEGIN;
  UPDATE orders SET status = 'shipped' WHERE id = 42;
  INSERT INTO outbox (topic, key, payload)
  VALUES ('orders.events', '42', '{"type":"OrderShipped","order_id":42}'::jsonb);
COMMIT;
-- Si el COMMIT falla → ambos rollbackean
-- Si el COMMIT exitoso → el outbox poller publicará el mensaje

-- El poller (puede ser pg_cron o un worker externo):
BEGIN;
  SELECT id, topic, key, payload
  FROM outbox WHERE sent_at IS NULL
  ORDER BY created_at LIMIT 100
  FOR UPDATE SKIP LOCKED;
  -- Publicar a Kafka aquí...
  UPDATE outbox SET sent_at = NOW() WHERE id IN (...);
COMMIT;
\end{lstlisting}

\section[Replicación Lógica para CDC]{\emj{🔄}~Replicación Lógica para CDC}

CDC (Change Data Capture) convierte el WAL de PostgreSQL en un stream de eventos de negocio.
En lugar de que la aplicación publique eventos explícitamente (Outbox), CDC los captura directamente de la BD observando los cambios en las tablas.
Esto es útil para integrar sistemas legados que no pueden modificarse, sincronizar datos a un data warehouse, o alimentar un pipeline de Kafka sin cambiar el código de la aplicación.

La replicación lógica nativa de PostgreSQL permite que otra base de datos se suscriba a los cambios en tablas específicas.
Herramientas como Debezium van más lejos: leen el WAL lógico y publican cada INSERT/UPDATE/DELETE como mensajes en Kafka, con la información completa de antes y después del cambio.

Change Data Capture (CDC) captura cambios de la BD como un stream de eventos.
La replicación lógica de PostgreSQL es una implementación nativa de CDC.

\begin{lstlisting}[caption={Replicación lógica como CDC}]
-- Requiere wal_level = logical en postgresql.conf

-- Crear publication (qué cambios capturar):
CREATE PUBLICATION orders_pub FOR TABLE orders, items;
-- O capturar todo:
CREATE PUBLICATION all_tables FOR ALL TABLES;

-- Crear subscription en otra base (que recibirá los cambios):
CREATE SUBSCRIPTION orders_sub
  CONNECTION 'host=source_db port=5432 dbname=prod user=replicator password=secret'
  PUBLICATION orders_pub;

-- Ver estado de la replicación lógica:
SELECT * FROM pg_stat_subscription;
SELECT * FROM pg_replication_slots;

-- Herramientas de CDC más maduras:
-- Debezium: captura WAL lógico → Kafka. Soporta PG, MySQL, MongoDB
-- pglogical: extensión para replicación lógica bidireccional entre PG
\end{lstlisting}

\begin{entrevistabox}
\begin{enumerate}
  \item ¿Cuál es la diferencia fundamental entre CRUD tradicional y Event Sourcing?
  \item ¿Qué es un aggregate en el contexto de Event Sourcing?
  \item Explica el problema que resuelve el Outbox Pattern. ¿Qué garantía de delivery ofrece?
  \item ¿Qué es eventual consistency en CQRS? ¿Cómo lo manejas en la UI?
  \item ¿Cuándo elegiría Event Sourcing sobre un modelo relacional clásico?
  \item ¿Qué es CDC (Change Data Capture) y cómo lo implementarías en PostgreSQL?
\end{enumerate}
\end{entrevistabox}
